% --- Chapter 1: Introduction ---
\chapter{INTRODUCTION}
% --- Project Information Table ---
\begin{table}[H]
\centering
\renewcommand{\arraystretch}{1.5}
\begin{tabular}{|l|l|}
\hline
\rowcolor[HTML]{EFEFEF} 
\textbf{Field} & \textbf{Details} \\ \hline
\textbf{Project Title} & Event Management \& Online Ticketing System \\ \hline
\textbf{Project ID} & ITPJ2602 \\ \hline
\textbf{Company Name} & Dề Dê Startup \\ \hline
\textbf{Project Manager} & Lê Nguyễn Hoàng Phúc - 23521198 \\ \hline
\textbf{Author} & Odoo Master Company \\ \hline
\textbf{Project Begin Date} & March 20, 2026 \\ \hline
\textbf{Client Date Submitted} & March 12, 2026 \\ \hline
\textbf{Current Version} & 1.0 \\ \hline
\textbf{End Date} & August 20, 2026 (Estimated 5-month duration) \\ \hline
\end{tabular}
\caption{General Project's Information table}
\end{table}

\textbf{Statement of Work Summary:} \\
The project aims to develop a comprehensive digital ecosystem consisting of a Web Application and a Mobile Module designed to automate the ticketing process, secure event entry via QR scanning, and provide deep business analytics for De De Startup. The system is engineered to handle high-traffic bursts and eliminate manual processing risks.

\section{Background}
De De Startup is a rapidly growing event organizer managing an average of 15–20 high-profile events annually, including concerts, seminars, and workshops. Currently, the organization relies on manual processes via Facebook and direct bank transfers, which has led to critical operational failures:
\begin{itemize}
    \item Proliferation of counterfeit tickets due to lack of secure verification.
    \item Frequent system crashes during "Golden Hour" sales windows.
    \item Inability to manage seat inventory effectively.
    \item Significant loss of revenue and lack of customer behavioral data for marketing.
\end{itemize}

\section{Scope of Work}
The project entails the end-to-end development of a scalable ticketing platform. Key functional domains include:
\begin{itemize}
    \item \textbf{Web Platform:} Event creation, dynamic seating map management, simulated online payment integration, and comprehensive financial reporting.
    \item \textbf{Mobile Module:} A specialized application for ticket controllers to scan QR codes with short-term offline capability.
    \item \textbf{Technical Constraints:} 1.5 Billion VNĐ budget, 5-member core team, and a 5-month timeline.
    \item \textbf{Delivery Method:} The project will be executed using an \textbf{Incremental and Iterative} approach, following the \textbf{Scrum} Project Management methodology to ensure flexibility and continuous value delivery.
\end{itemize}

\section{Objectives}
The following table outlines the primary objectives of the project and the strategic rationale behind them:

\begin{table}[H]
\centering
\renewcommand{\arraystretch}{1.5}
\begin{tabular}{|p{0.35\textwidth}|p{0.6\textwidth}|}
\hline
\rowcolor[HTML]{EFEFEF} 
\textbf{Objective} & \textbf{Rationale} \\ \hline
Handle 10,000 concurrent users (simulated). & To ensure system stability and prevent crashes during high-demand "Golden Hour" ticket releases. \\ \hline
System response time < 2 seconds. & To provide a seamless user experience and minimize customer drop-off rates during checkout. \\ \hline
Entry verification < 5 seconds per person. & To expedite the check-in process at event gates and prevent crowd congestion. \\ \hline
Increase ticket revenue by 20\%. & To recoup investment through automated sales, better inventory management, and data-driven marketing. \\ \hline
Zero duplicate QR codes and 98\%+ payment success rate. & To eliminate fraud, secure financial transactions, and maintain brand reputation. \\ \hline
Implement basic anti-fraud and offline scanning. & To ensure operational continuity at venues with unstable internet and protect against bot attacks. \\ \hline
\end{tabular}
\end{table}